\section[Selected Green entries]{Selected Green entries from the same cluster records}
\label{sec:op3-cluster-green}

\paragraph{Status.}
The query identities below are \statusproved{} as a draft awaiting
independent review and are implemented in an exact audit. Replacing the
coarse solve of \cref{thm:op3-local-cycle-rank} by a complete maintained
inverse algorithm remains \statusopen{}. A block inverse formula alone
does not pay for changing retained vertices, selecting covariance entries
or updating every stored matrix word.

\begin{lemma}[Conditional variance and selected covariance; \statusproved]
\label{lem:op3-cluster-conditional-green}
Let a valid tree-piece hierarchy have height $h$, root ports $P_B$ and
interior $I_B$. For a named vertex $i$, one incident-edge parent walk returns
its conditional affine response
\[
 w_i=z_i+\boldsymbol a_i^{\mathsf T}\boldsymbol w_{P_B}
\]
and conditional variance
$g_{ii}=(\boldsymbol M_{I_BI_B}^{-1})_{ii}$, extended by zero when $i$
is a port, in $O(1+h)$ word work and constant live arithmetic state.
Optionally retain $O(1+h)$ scalar noise coefficients from that walk.
For two vertices in the same piece, their retained traces give
$g_{ij}=(\boldsymbol M_{I_BI_B}^{-1})_{ij}$, again extended by zero at
ports, in $O((1+h)\log(2+h))$ comparison-dictionary work.
Thus the height-bounded source hierarchy gives logarithmic named variance
work and at most $O(\log^2(2+n))$ selected covariance work.
\end{lemma}
\begin{proof}
At a compress or forget operation, let $m$ be the eliminated scalar and
$\delta_m>0$ its pivot. Completing its square factors the conditional
quadratic into
\[
 \tfrac12\delta_m
 \bigl(w_m-\boldsymbol\eta_m^{\mathsf T}\boldsymbol w_{\rm parent}
              -\sigma_m\bigr)^2
\]
plus the parent Schur quadratic. Raking joins disjoint interiors and
introduces no eliminated scalar. Applying this identity recursively gives
a triangular factorization with one independent scalar residual
$\xi_m$ of variance $1/\delta_m$ per eliminated interior vertex. This is
an algebraic covariance interpretation of the inverse factorization;
the algorithm samples no random variables.

Follow the affine pullback of \cref{lem:op3-cluster-point-query}. Whenever
the current row has coefficient $c_m$ on an eliminated port, add
$c_m^2/\delta_m$ to a variance accumulator and, optionally, save
$(m,c_m,1/\delta_m)$ keyed by its current elimination record. The rest of
the row pulls back to the parent ports as before. Previously collected
residual coefficients do not change, since the factorization makes those
residuals independent of the remaining coordinates. At the root,
\[
 w_i=z_i+\boldsymbol a_i^{\mathsf T}\boldsymbol w_{P_B}
          +\sum_m c_{i,m}\xi_m,
 \qquad
 g_{ij}=\sum_m\frac{c_{i,m}c_{j,m}}{\delta_m}.
\]
Only eliminated records on the query's ancestor path can occur. Each
structural callback has at most two helper records, so there are $O(1+h)$
terms. A compress pivot is the sum of the two children's shared diagonal
entries minus the original degree; a forget pivot is the child's eliminated
diagonal entry. Both are read or computed in constant work from existing
records. Matching two saved traces in a comparison dictionary gives the
claimed covariance cost. A root port has its basis response and an empty
noise trace, as required. Different pieces have conditionally independent
interiors given all coarse ports; their conditional cross term is zero by
the partition, not merely because two arbitrary record identifiers differ.
\end{proof}

\begin{lemma}[Port promotion and an original-vertex border; \statusproved]
\label{lem:op3-green-promotion-border}
Suppose the full current active inverse restricted to retained ports is
$\boldsymbol J=(\boldsymbol M_{UU}^{-1})_{PP}$.
Extend each piece response row $\boldsymbol a_i$ by zero to all of $P$.
For old active vertices $i,j$, its full Green entry is
\begin{equation}
 (\boldsymbol M_{UU}^{-1})_{ij}
   =g_{ij}+\boldsymbol a_i^{\mathsf T}\boldsymbol J\boldsymbol a_j,
 \label{eq:op3-conditional-full-green}
\end{equation}
where $g_{ij}=0$ across different piece interiors and at retained ports.
In particular, promoting old interior vertices to retained coordinates
augments $\boldsymbol J$ with these entries without solving a new full
active system.

If a new candidate $j$ has active-parent indicator $\boldsymbol s$ on $U$,
put $\boldsymbol g=\boldsymbol M_{UU}^{-1}\boldsymbol s$ and
$\delta=d_j-\gamma^2\boldsymbol s^{\mathsf T}\boldsymbol g$.
Then $\delta\geq\bar\alpha d_j>0$ and the enlarged inverse is
\begin{equation}
 \begin{pmatrix}
  \boldsymbol M_{UU}^{-1}+\gamma^2\boldsymbol g\boldsymbol g^{\mathsf T}/\delta
       &\gamma\boldsymbol g/\delta\\
  \gamma\boldsymbol g^{\mathsf T}/\delta&1/\delta
 \end{pmatrix}.
 \label{eq:op3-original-border-inverse}
\end{equation}
For the original physical face values,
$u_j=(\gamma\boldsymbol s^{\mathsf T}\boldsymbol u_U-\lambda d_j)/\delta$
and the old values increase by $\gamma\boldsymbol g u_j$.
\end{lemma}
\begin{proof}
The inverse factorization underlying
\cref{lem:op3-cluster-conditional-green} separates conditional interior
residuals from the retained coordinates, proving
\eqref{eq:op3-conditional-full-green}. Equivalently, it is the usual block
inverse of a matrix partitioned into interiors and ports. Each response
row has at most two nonzero retained coefficients.

The enlarged original matrix has blocks
$\boldsymbol M_{UU}$, $-\gamma\boldsymbol s$ and $d_j$.
Direct block elimination proves \eqref{eq:op3-original-border-inverse}
and the physical value update. For any principal active set $W$ and vector
$\boldsymbol y$, the original full degrees give
\[
 \boldsymbol y^{\mathsf T}\boldsymbol M_{WW}\boldsymbol y
 \geq \sum_{i\in W}(d_i-\gamma d_{G[W]}(i))y_i^2
 \geq\bar\alpha\sum_{i\in W}d_i y_i^2.
\]
Minimizing the enlarged quadratic over old coordinates with the new
coordinate fixed to one gives $\delta\geq\bar\alpha d_j$.
The identities are exact; this lower bound by itself is not a floating-point
stability theorem for a long sequence of inverse updates.
\end{proof}

\paragraph{Precise remaining algorithmic obligation.}
For an ordinary one-parent admission at $i$, a named trace supplies
$g_{ii}$ and the at most two nonzero coefficients of $\boldsymbol a_i$.
Then the selected old-port Green vector is
$\boldsymbol J\boldsymbol a_i$ and the full parent variance is
$g_{ii}+\boldsymbol a_i^{\mathsf T}\boldsymbol J\boldsymbol a_i$.
Equation \eqref{eq:op3-original-border-inverse} suggests a rank-one update
of the stored port inverse in $O(|P|^2)$ work, with every inverse entry
written and charged.

At a cycle birth, first promote the old active parents and newly needed
Steiner junctions using selected covariances, then border the new vertex
with all its active edges. At most one parent, the insertion-tree parent,
may be a transient port; restricting the updated inverse afterward to the
permanent retained set preserves the original partition's port bound.
Retaining every promoted marker is an alternative with a larger constant.
This step must
be reconciled with the paid repartition and with queries from the old
current piece versions. The full update implementation, exact comparisons
and total promotion/copy bound are the next falsifiable target in
\texttt{COARSE\_INVERSE\_UPDATE\_PROBE.md}. No improved end-to-end exponent
is asserted from these identities alone.

\paragraph{Measured.}
\texttt{cluster\_green\_queries.py} compares affine rows, conditional
variances and all covariance pairs against independent exact conditional
inverse matrices. It includes one-port bases, all compress/rake/forget
helpers, arbitrary physical seeds, shifted source-interior clusters and
retained payload versions. Whole-hierarchy indexing and dense inverse
columns are separate reference costs. Reproducible counts and hashes are
recorded in \texttt{CLUSTER\_GREEN\_QUERY\_AUDIT.json}. The sweep includes
256 canonical faces, 2,664 independent conditional inverse matrices,
29,408 affine/variance checks and 123,844 covariance pairs, with no mismatch.
